\section{Actual Archives and the Size of the Benefit}\label{sec:exp}
We distinguish an analysis of genuinely incomplete records from controlled
comparisons where complete records allow exact evaluation. The new tasks,
splits and summaries were specified before fitting their models.

\paragraph{A frozen comparison on preassigned PISA forms.}
The U.S. public-use student file contains 4,978 students in 162 schools
\citep{pisaus2014}. We use the common report of grade repetition in grades
1--6 as the binary target and three binary response indicators from each
of three rotated question sets. The estimand concerns reported responses;
asked-but-unanswered items count as nonendorsements, while items absent
from a student's assigned form remain unobserved. The 226 records without
a valid common target response are excluded, defining the respondent
population being compared. The original instrument's limit is 30 minutes;
six of nine refers only to our selected predictors.

Forty schools train the old ensemble, another forty train the new one, and
the remaining 82 supply 2,441 audit records. Each ensemble averages three
block-specific forests, fitted using only students who received that
block. The contrast depends on all nine predictors. Every pair of blocks
occurs in a form, so both the accuracy and Brier differences decompose
into observable functions (Appendix~\ref{app:pisa}). We freeze the models
and form-score tables before examining the audit outcomes. Calibration
uses incomplete training records; there are no complete reference records
with which to perform a direct audit.

\begin{table}[t]
\centering
\caption{PISA secondary analysis on 2,441 held-out students. Brier improvement and 95\% intervals are multiplied by $10^3$; positive means the newer rule has smaller loss. Survey intervals use the 80 supplied Fay replicates. MI intervals also use Rubin combination under a working imputation model.}
\label{tab:pisa}
\small
\begin{tabular}{lrr}
\toprule
Method & Estimate & 95\% interval \\
\midrule
Form score & 1.13 & $[-2.28, 4.53]$ \\
Block/pair HT & 1.24 & $[-2.09, 4.57]$ \\
Logistic MI & 1.24 & $[-2.12, 4.60]$ \\
\bottomrule
\end{tabular}
\end{table}


Table~\ref{tab:pisa} compares optimized scores, a simple inverse-probability
block/pair score, and twenty logistic chained-equation imputations with
Rubin combination \citep{vanbuuren2011}. Survey-weighted intervals use the
80 supplied Fay replicates \citep{kastberg2014}. All three Brier intervals
include zero: these data do not establish an improvement in probability
quality. The methods give similar estimates and widths; optimization is
not automatically better after fitting. The randomized-accuracy gain is
about 4.1 percentage points, a different target. Both rules threshold to
the negative class everywhere on this cube, so their deterministic
accuracy difference is zero.

The design interpretation assumes invariant reported answers across forms
and randomized assignment retained in the common-target respondent
population. Nine duplicated-item diagnostics find two unadjusted
$p<0.05$ differences, neither surviving Bonferroni correction. This check
does not verify invariance. MI adds a working imputation model. Because
complete-data truth and a branching counterfactual are unavailable, this
analysis measures neither coverage nor an adaptivity ratio.

\paragraph{Independent forest updates under a binding cap.}
On WDBC, Spambase and Ionosphere, twenty fixed splits each fit independent
old and new random forests, using disjoint training records
\citep{wolberg1993,spambase1999,ionosphere1989}. Each forest has twenty
depth-three trees. Nested sets of 3, 6, 9 and 12 training-binarized features
give 240 tasks, all with cap $b=3$. Probability averages remain degree at
most three, even when their active feature set is larger. This statement
does not cover thresholding the ensemble. Existing model-training records
calibrate the score tables, and held-out records evaluate their exact
conditional variances.

The fixed menu contains every feature triple. The branching menu contains
those forms and all forty component trees, with jointly optimized leaf
scores. It can choose a component tree and follow its path, but cannot
evaluate the whole forest in three queries. Thus $\Var(Yd)$ is a lower
bound, not the denominator of the reported branching comparison.

\begin{table}[t]
\centering
\caption{Independent forests, cap 3: 60 binding tasks per dataset. Median fitted fixed/menu-branch variance ratio and oracle menu ratio (numerical lower endpoint). The branch menu includes every fixed triple and all 40 component trees. Neither column uses the unattainable full-forest target as its branching score.}
\label{tab:forest}
\small
\begin{tabular}{lrrr}
\toprule
Data & Fitted & Oracle menu & Fitted $>1.5$ \\
\midrule
WDBC & 1.085 & 1.090 & 0/60 \\
Spambase & 1.029 & 1.031 & 0/60 \\
Ionosphere & 1.117 & 1.100 & 0/60 \\
\bottomrule
\end{tabular}
\end{table}


All 180 tasks with 6--12 candidate features have more than three active
coordinates. Their fitted variance ratios have medians 1.03--1.12 across
datasets; none of the 240 tasks exceeds 1.5. Figure~\ref{fig:forestbinding}
reports every task and its distribution. Oracle optimization on the same
menus gives similar modest gaps. Its fixed-form lower bound divided by a
feasible branching variance lower-bounds the unrestricted intrinsic ratio;
its upper endpoint only concerns the supplied branch menu. More elaborate
branching strategies can still improve on that menu.

This also clarifies the smaller independent-tree experiment. Among its
70 probability-output comparisons, 60 already have active feature count
above the cap. Their population ratios have medians 1.353 (SUSHI) and
1.295 (SAPA), while global coefficients remain near equality. For
threshold outputs, 52 of 58 nonzero comparisons are nonbinding. Structural
ties explain most threshold results, but do not explain away the modest
probability-output gaps (Appendix~\ref{app:binding}).

\paragraph{Controlled gaps, transfer and operational costs.}
Depth-six CART/root refinements isolate a contrast that branching can
compute in six queries. Exhaustive lower bounds over all fixed six-feature
forms certify median population ratios at least 1.883 on SUSHI and 1.474
on SAPA. These are controlled stress tests, distinct from the independent
forest updates. The complete-calibration version is an efficiency
diagnostic: its 500 independent target-population records already suffice
for a direct audit at the chosen precision. Transferring calibration
across SAPA age groups avoids that interpretation. Brier and clipped
log-loss comparisons retain median fitted variance ratios 2.350 and 2.623,
with conditional Wald coverage near 95\%. A separate restricted-pilot
experiment charges all newly collected calibration records
(Appendices~\ref{app:transfer} and~\ref{app:em}).

If the cap can be relaxed, compare complete acquisition using its actual
label cost. With fixed/complete variance ratio $\rho>1$, fixed form size
$b$ and complete size $m$, their continuous cost crossover is
\begin{equation}
 c_Y^*=\frac{m-b\rho}{\rho-1}.\label{eq:costcross}
\end{equation}
Median crossovers on the controlled tasks are 3.16 and 3.74 feature-cost
units. At $c_Y=10$, complete records cost less on 45/45 SUSHI and 19/25
SAPA tasks. Figure~\ref{fig:costchoice} makes this conditional choice
explicit. When the cap is immutable, complete acquisition is inadmissible.
For nonlinear updates, the bias certificate must be checked first:
WDBC GBDT/MLP budget curves reject overly tight caps regardless of sample
size (Appendix~\ref{app:budgetcurve}).
